\subsection{Further quantitative statements}

The rest of the paper has three parts.  The first proves the uniform-bound
criterion and its specialization dichotomy.  The second makes a single
Weil detector effective, transports detectors between lattices, and derives
degree bounds from specified geometric inputs.  The third studies the
remaining existence questions: it gives constructions with precisely stated
deformation gaps, proves obstructions for particular constructions, and
develops a separate detector criterion for Mumford fourfolds.

For a polarized abelian variety of Weil type, an imaginary quadratic field $K$ acts with signature $(n,n)$ on a variety of dimension $2n$. Its rational plane of Weil classes is
$W_K=\bigwedge_K^{2n}H^1(A,\Q)$. Let $\delta_W$ be the minimum degree of an irreducible subvariety with nonzero projection to $W_K$. Theorem~\ref{thm:extraction} proves, for every integral polarization compatible with the action, every nonzero $t\in W_K$, every $\delta\in\Z_{>0}$ with $\delta t$ integral, and every endomorphism $a\in K$ with $a/\bar a$ not a root of unity,
\[
 \delta_W\leq c_t\leq
 \left\lceil\frac{2\,16^n}{n!}\right\rceil
 \delta\lceil\sqrt T\rceil\Nm(a)^{n(2n+1)}\delta_W^2,
 \qquad T=(-1)^n\int t^2.
\]
Compatibility means that the Rosati involution restricts to conjugation on $K$. Infinite values have their usual extended-real interpretation. Thus one detector gives the whole Weil plane at quadratic cost.

Theorem~\ref{thm:small-detectors}, using the cycles of Schoen and van Geemen, gives detector degrees $48$ for even principal Gaussian fourfolds, $192$ for odd principal Gaussian fourfolds, and $36$ for principal Eisenstein fourfolds, all with full integral maximal-order action. In the Eisenstein case $3(W_K\cap H^4(A,\Z))$ is algebraic and
\[
 c_t\leq\max\{3\delta,36\lfloor\delta\sqrt{2T}/3\rfloor\}
 \leq17\delta\lceil\sqrt T\rceil.
\]
The even and odd Gaussian lattices, the integral normalization and the effective cycles are specified in Section~\ref{sec:small}.

Theorem~\ref{thm:stabilization}, after Schoen, says that bounded detectors in every connected split Weil family of dimension $2n+2$ over a fixed imaginary quadratic field imply bounded detectors in every Weil family of dimension $2n$ over that field, for every determinant. No class-number restriction is needed. If $\mathcal O_K=\Z[\tau]$ has class number one, acts integrally, and $\Delta=\chi(L)$, Theorem~\ref{thm:explicit-stabilization} gives
\[
 \delta_W(A,L)\leq\frac{\max\{1,\Nm(\tau)/\Delta\}}{n+1}D_{n+1}(\Delta^2),
\]
where $D_{n+1}(x)$ bounds detectors for all split Weil $(2n+2)$-folds of polarization product $x$. The premise includes every polarization type of that product. The numerator of the displayed factor is one for the Gaussian and Eisenstein fields.

Schoen's cyclic cycles and the dominance theorems of Schoen and Koike give the explicit sixfold estimates in Theorem~\ref{thm:sixfold-suppliers}: $746496$ for the even principal Gaussian lattice, $5971968$ for the odd one, and $30375$ for the principal Eisenstein lattice. Theorem~\ref{thm:principal-quotient} gives a principal hyperbolic quotient over every imaginary quadratic field. Pullback of detectors costs $t_K^nf^n\Delta$, where $\operatorname{Tr}(\mathcal O_K)=t_K\Z$, $f$ is the conductor of the acting order and $\Delta$ is the polarization volume. The ideal classes can also be removed at the additional cost $m_K^n$, where $m_K\le2\sqrt{|\operatorname{disc}K|}/\pi$. Thus one principal family with lattice $U^n$ suffices for each field and dimension; for class number one the additional factor is one. Consequently, for a compatible polarized Weil fourfold of any determinant with integral action by $\mathcal O_f=\Z+f\mathcal O_K$ and $\Delta=\chi(L)$, Theorem~\ref{thm:all-determinant-fourfolds} gives
\[
 \delta_W\le1990656 f^6\Delta^2\quad(K=\Q(i)),
 \qquad\delta_W\le10125 f^6\Delta^2\quad(K=\Q(\omega)).
\]
The proof keeps the source cycles effective and uses published geometric inputs.

Markman \cite[Theorem~1.2]{MarkmanICM} proves algebraicity for split Weil sixfolds over every imaginary quadratic field, and for all Weil fourfolds; the published account gives a proof sketch and points to the detailed preprint. The field restriction to nine class-number-one fields belongs to particular numerical estimates, not to this qualitative theorem. Assuming Perry's unrefereed equivariant semiregularity theorem, Theorem~\ref{thm:all-field-sixfolds} gives an explicit bound for every imaginary quadratic field: if $K=\Q(\sqrt{-d})$, $d$ is positive squarefree, $\mathcal O_f$ acts integrally, and the compatible polarization has volume $\Delta$, then every split Weil sixfold satisfies
\[
 \delta_W<10^{179}d^{53}f^3\Delta.
\]
For class number one the coefficient improves to $10^{178}d^{51}$. The sharper constant is an explicitly evaluated Chern-class degree. Split eightfolds and higher dimensions require further arguments.

The $G_{34}$ family in Severn's unrefereed preprint gives a different input in dimension eight. Theorem~\ref{thm:eisenstein-eightfolds} writes out an explicit integer $B_0$ and proves, from that construction, the bounds
\[
 \begin{aligned}
 \delta_W&\le B_0f^4\Delta
 &&\text{for split Eisenstein eightfolds},\\
 \delta_W&\le\frac{B_0}{4}\frac{f^9\Delta^2}{b}
 &&\text{for all Eisenstein sixfolds},
 \end{aligned}
\]
where $b$ is the index of maximal-order saturation. Schoen's cyclic-cover cycles, Severn's exceptional-divisor correction, Brauer induction and Milne's invariant theory carry this proof. The degree calculation keeps integral orbit multiplicities and then uses the principal quotient to reduce the transport cost.

Theorem~\ref{thm:ghosts} applies to any fixed smooth projective variety over a number field with an ample line bundle and any rational Betti class. Its minimum numerator degree at good reductions is eventually the characteristic-zero minimum if that is finite; otherwise it exceeds every fixed bound outside finitely many places. The multiplier and auxiliary prime may vary. The same dichotomy holds for the cost that also bounds the multiplier.

These results concern algebraic representatives, not merely Hodge-theoretic tensors or motivated classes. Andr\'e's deformation theorem and his theorem that abelian Hodge classes are motivated \cite[Theorems~0.5 and 0.6.2, pp.~8--9]{Andre96} establish powerful propagation results in the motivated category. His CM reduction \cite[Lemma~6.3.2, p.~32]{Andre96} uses split Weil classes over general CM fields. Neither assertion supplies the uniform algebraic cycles required here, and the present criterion does not reduce all abelian Hodge classes to imaginary-quadratic Weil classes.

Two restricted eightfold constructions admit precise tests. Theorem~\ref{thm:theta-rank} shows that, for $V$ with basis indexed by $\mathbf F_2^4$, every nonzero alternating four-form invariant under translations and character signs induces a pairing on $\bigwedge^2V$ of rank at least $24$, and that $24$ occurs. Corollary~\ref{cor:all-eightfold-quadrics} identifies this $35$-dimensional space with the entire quadratic ideal of the minus projection on a nonempty open subset of the even Gaussian eightfold family. Thus no quadratic relation there can provide the rank-$10$ quadric needed by the direct ruling-class extension of van Geemen's argument. Theorem~\ref{thm:localized-support} proves a different obstruction: a perfect complex supported on translates of abelian subvarieties meeting in codimension at least two cannot deform in a direction with nonzero normal class on a component of nonzero generic virtual multiplicity, even with an infinitesimal gerbe twist. On norm-one scalar graphs in $E_i^8$ this excludes all six skew normal directions, including for arbitrary extensions and thickenings on those supports. The general eightfold existence problem remains open.


Two other results give exact tests for possible eightfold constructions. Theorem~\ref{thm:pure-weil-ext} applies to a perfect complex $E$ on a balanced Weil $2n$-fold, $n\ge2$, with a rational algebraic $B$-field and an exactly pure nonzero normalized Weil character. For a coherent finite translation--tensor action and an equivariant structure, its ordinary semiregularity trace on invariant obstructions is injective if and only if
\[
 \dim\operatorname{Ext}^2(E,E)^G=2n(2n-1).
\]
On an eightfold this number is $56$. Proposition~\ref{prop:sixteen-line-character} constructs an integral sixteen-term character on the principal Gaussian family with precisely that normalized Weil character, but Proposition~\ref{prop:sixteen-line-gluing} proves that every indicated two-level gluing has invariant second self-Ext dimension $448$. Theorem~\ref{thm:all-point-products} excludes every external product of point-ideal characters, for every partition of dimension eight, and all their rational Fourier--Mukai orbits as sources of a nonzero full-family Weil component.

For finite covers of curves, Theorem~\ref{thm:finite-cover-rank} says that a faithful irreducible coefficient representation with nonrational character, eight-dimensional parabolic multiplicity of type $(4,4)$, and branch moving ranks at least two needs coefficient rank at least twenty-six to have a sixteen-dimensional period image. Theorem~\ref{thm:gaussian-reflections} excludes dominance when a Gaussian representation with exact character field contains a reflection. Hence the rank bound applies to every dominant source of that specified Gaussian kind. Theorem~\ref{thm:scalar-standard-all} further excludes scalar-standard coefficients in every degree: their period variations factor through rank-one local systems on curves with moduli dimension at most fifteen. These statements restrict construction methods and prove no additional case of algebraicity.


Theorem~\ref{thm:simple-hodge-bundle} gives simple vector bundles on
every CM elliptic product $E^{2n}$, for every imaginary quadratic
field, whose Chern character is a nonzero prescribed primitive
middle class plus an explicit polynomial in the polarization.
For a Weil class and $n\ge3$, the Hochschild action has rank
$7n^2-2n$, equal to $104$ on an eightfold; the general rank formula
is \cite[Theorem~A6]{BBK3Sixfolds}.
The construction uses positive multiplicities and ordinary extensions.
Theorem~\ref{thm:tree-extension-obstruction} proves that its specified
tree of permutation orbits cannot deform in every Weil direction,
even after its filtration is forgotten.  Additional gluing data and
the sufficient second-Ext equality remain to be found.
Theorem~\ref{thm:sign-cube-obstruction} extends the obstruction to a
specified family of coordinate-sign networks with many cycles.
Theorem~\ref{thm:gaussian-positive-trace-network} gives a different,
explicit Gaussian reflection network whose positive multiplicities
have a nonzero Weil character and solve every matrix-valued diagonal
trace equation.  The remaining off-diagonal extension equations and
strong semiregularity are not established, so this construction does
not yet give a deforming detector.

For a polarized Hodge-generic Mumford fourfold $A$, the exceptional
codimension-two classes on $A^2$ form a plane identified with the
trace-zero part of its defining totally real cubic field. Theorem
\ref{thm:mumford-cubic} proves
$\delta_{\mathrm{exc}}\leq c_t\leq C_t\delta_{\mathrm{exc}}^3$
for each fixed nonzero class $t$ in that plane, where
$\delta_{\mathrm{exc}}$ is the least degree of an irreducible
codimension-two detector for the product polarization. The constants
are uniform in a fixed polarized family with global $t$.
Corollary~\ref{cor:mumford-family} therefore reduces algebraicity
throughout that family to bounded detectors on infinitely many CM
fibres. These exceptional classes are not Weil classes; the theorem
does not construct their first detector.

When $A$ is the Jacobian of one of the genus-four curves of Bouchet--Hanselman--Pieper--Schiavone, Theorem~\ref{thm:mumford-surface} gives a further sufficient condition. Twist the polarization on its rank-nine K3-type Hodge structure by a totally positive $\lambda\in F\setminus\Q^\times F^{\times2}$, and realize it by a projective K3 surface $S_\lambda$. One algebraic surface correspondence from $S_\lambda$ to $C\times C$ with nonzero transcendental action makes the entire exceptional plane algebraic. The extraction uses a nonconstant weighted Casimir tensor. A dominant rational map from $C\times C$ to this chosen $S_\lambda$ cannot provide it. The theorem refines the conditional K3 approach in their Section~5.2; it does not construct the first correspondence.

A different conditional reduction is Theorem~\ref{thm:mumford-coniveau}. Suppose the quaternion algebra over $F$ is split by $F(\sqrt{-3})$ and has split corestriction to $\Q$. On a common finite cover of the Mumford curve there is a split Weil sixfold $B$ and an explicitly algebraic projector onto a rank-sixteen substructure $P\subset H^3(B,\Q)$ with $P(1)\simeq H^1(A,\Q)^{\oplus2}$. At a Hodge-generic point, one nonzero geometrically coniveau-one part of $P$ implies the rational Hodge conjecture for every power of $A$, including the exceptional plane on $A^2$. The known Weil cycles on $B$ give $P$ as a summand of the middle cohomology of supporting threefolds; the additional divisor support remains an open generalized Hodge question.
